\documentclass[Numbered]{sn-jnl}
\usepackage[numbers]{natbib}
\usepackage{amsmath,amssymb,booktabs,tabularx,graphicx}
\begin{document}
\title{Supplementary Information: Situation engineering is a missing layer for reliable artificial intelligence}
\author{Geunsik Lim}\maketitle

\section*{Supplementary Note 1: Benchmark schema}
Each JSONL item contains id, category, query time, question, claims, required variables, gold action and gold answer. Claims contain text and auditable event metadata.

\section*{Supplementary Note 2: Complete aggregate results}
\begin{table}[h]\centering\begin{tabular}{lrrr}\toprule Method & Accuracy & SCHR & ECE\\\midrule
SCQA & 1.000 & 0.000 & 0.053\\
Latest-mention & 0.857 & 0.143 & 0.183\\
Recency-RAG & 0.714 & 0.286 & 0.144\\
Lexical-RAG & 0.700 & 0.300 & 0.094\\
\bottomrule\end{tabular}\end{table}

SCHR is the complement of exact accuracy in this single-answer controlled
setting and is not interpreted as an independent outcome. ECE is computed from
rule-assigned confidence with small seeded perturbations; it is retained only as
a software regression smoke test and is not evidence of neural calibration.

\section*{Supplementary Note 3: Category results}
\begin{tabular}{lrrrr}
\toprule
category & Latest-mention & Lexical-RAG & Recency-RAG & SCQA \\
\midrule
counterfactual & 1.000 & 1.000 & 1.000 & 1.000 \\
hidden\_premise & 0.000 & 0.000 & 0.000 & 1.000 \\
modality & 1.000 & 1.000 & 1.000 & 1.000 \\
observer & 1.000 & 0.530 & 0.000 & 1.000 \\
scope & 1.000 & 1.000 & 1.000 & 1.000 \\
source\_conflict & 1.000 & 1.000 & 1.000 & 1.000 \\
temporal & 1.000 & 0.372 & 1.000 & 1.000 \\
\bottomrule
\end{tabular}


\section*{Supplementary Note 4: Ablations}
\begin{tabular}{lrr}
\toprule
ablation & accuracy & errors \\
\midrule
none & 1.000 & 0 \\
time & 0.935 & 273 \\
source & 1.000 & 0 \\
modality & 0.857 & 600 \\
scope & 0.857 & 600 \\
observer & 0.857 & 600 \\
world & 1.000 & 0 \\
missing & 0.857 & 600 \\
\bottomrule
\end{tabular}


\section*{Supplementary Note 5: Reproduction}
From the repository root, run
\texttt{PYTHONPATH=code python code/run\_experiments.py}. This regenerates
gold-, predicted- and corrupted-state outputs. Run
\texttt{PYTHONPATH=code pytest -q code/tests} for implementation tests.

\section*{Supplementary Note 6: Situation-engineering design checklist}
\begin{table}[h]
\centering
\caption{Minimum reporting checklist for situation-engineered AI systems.}
\begin{tabularx}{\linewidth}{p{3.2cm}X}
\toprule
Component & Questions that should be answered \\
\midrule
Situation schema & Which actors, events, valid times, scopes, modalities, observers, worlds and uncertainty variables are represented? \\
Sensing & Which variables are supplied, retrieved or inferred? How is extraction accuracy measured independently of final answers? \\
Assembly and update & How are event coreference, supersession, conflict, validity intervals and belief revision handled? \\
Policy & Under which conditions does the system answer, retrieve, clarify, abstain or act? \\
Calibration & Does state uncertainty propagate into answer confidence and selective risk? \\
Observability & Can an auditor recover the active state, provenance, alternatives and policy rationale for each output? \\
Ablation & Does removing a state variable create the predicted category-specific failure rather than only a generic score reduction? \\
External validity & Are naturally occurring, multilingual, multimodal, adversarial and distribution-shift settings evaluated? \\
Governance & Are sensitive state variables minimized, access-controlled, correctable and audited for disparate effects? \\
\bottomrule
\end{tabularx}
\end{table}

\section*{Supplementary Note 7: Situation-engineering maturity model}
\begin{table}[h]
\centering
\caption{Proposed SE0--SE4 maturity levels.}
\begin{tabularx}{\linewidth}{p{1.0cm}p{3.4cm}X}
\toprule
Level & Capability & Falsifiable evidence \\
\midrule
SE0 & Implicit situation inference only & No external situation representation or state-specific evaluation \\
SE1 & Isolated situation tags & Measured extraction of individual variables such as time, location or modality \\
SE2 & Explicit dynamic state & Provenance-linked state, update rules and temporal or scope consistency tests \\
SE3 & Epistemic control & Missing-variable detection, competing hypotheses and answer/retrieve/clarify/abstain policy \\
SE4 & Adaptive audited situation intelligence & Learned multimodal sensing, multi-agent belief tracking, shift monitoring and independent state audits \\
\bottomrule
\end{tabularx}
\end{table}

\section*{Supplementary Note 8: Recommended future benchmark modules}
A natural-text extension should separately score event extraction, event coreference, temporal ordering, validity intervals, source and modality classification, scope resolution, observer belief, counterfactual frame selection, missing-variable detection, clarification utility, selective risk and claim-level situation consistency. Fixed-state policy tests and fixed-policy sensing tests are both needed to prevent final accuracy from obscuring the source of improvement.

\section*{Supplementary Note 9: Frontier-model evaluation status}
The repository includes a frozen temporal SituatedQA smoke-test sample and a
multi-provider evaluation harness. No frontier-model outputs are reported in
this version because the execution environment contained no provider API
credentials and no local model checkpoints. The harness aborts rather than
creating placeholder predictions. A publishable revision should run the full
licensed split on at least three model families, archive raw responses and
record exact model snapshots, dates, prompts, token budgets, costs and failed
calls. This distinction is essential: an executable protocol is not an
experimental result.


\section*{Supplementary Note 10: Independent annotation protocol}
The executable workflow creates separately randomized packets for at least three
annotators. Annotators label action, answer, temporal state, modality, scope,
source status, observer state and possible world without seeing method identity
or another annotator's labels. Files are frozen before agreement calculation
and adjudication. The scorer rejects missing cells and reports slot-level
Fleiss' $\kappa$ and unanimous agreement. Recruitment, ethics status,
compensation, qualifications, exclusions and adjudication changes must
accompany any reported human result. No such result is claimed here.

\section*{Supplementary Note 11: Interfaces among engineering disciplines}
\begin{table}[h]
\centering
\caption{Recommended typed interfaces between situation engineering and neighbouring components.}
\begin{tabularx}{\linewidth}{p{3.2cm}p{4.2cm}X}
\toprule
Neighbouring discipline & Input to the situation layer & Output or requirement from the situation layer \\
\midrule
Prompt engineering & Task, desired output, constraints and user intent & Required situation variables and clarification policy \\
Context / retrieval engineering & Candidate claims, passages, timestamps and source metadata & Missing-variable queries, applicability filters and freshness requirements \\
Memory engineering & Events, commitments, observations and provenance across sessions & Validity, supersession, observer/world tags and forgetting constraints \\
Tool engineering & Tool schema, preconditions, results, side effects and error states & Permitted calls, required confirmation and normalized state transition \\
Protocol engineering & Capability negotiation, authorization scopes and resource identifiers & Situation schema version, provenance envelope and uncertainty representation \\
Harness engineering & Workflow trace, checkpoints, permissions and recovery state & State invariants, action gate and answer/retrieve/clarify/abstain decision \\
Evaluation engineering & Test scenarios, judges, counterfactual perturbations and production traces & State-specific expected behaviour and decomposed failure labels \\
Observability engineering & Logs, lineage, timing and causal trace & Active-state snapshot, alternatives, unresolved variables and policy rationale \\
Serving engineering & Model availability, latency, token, memory and cost budgets & Minimum verification budget and graceful degradation policy \\
Security engineering & Trust zones, identities, privileges and threat signals & State-dependent authorization and response restrictions \\
\bottomrule
\end{tabularx}
\end{table}

\section*{Supplementary Note 12: Engineering anti-patterns}
The taxonomy enables failure analysis that is more precise than attributing every
problem to ``the prompt'' or ``the model''. Common anti-patterns include:
(i) context dumping, in which more documents are added without resolving which
state is active; (ii) memory accumulation, in which obsolete beliefs are retained
without validity or supersession metadata; (iii) tool proliferation, in which
capabilities are exposed without preconditions, typed errors or state effects;
(iv) harness opacity, in which retries and intermediate actions are not connected
to an auditable situation state; (v) evaluation collapse, in which only final
answer accuracy is measured and sensing, state update and policy failures are
indistinguishable; and (vi) serving-induced semantic degradation, in which cost
or latency optimizations silently reduce retrieval depth, truncate decisive
evidence or skip verification. A situation-engineered system should make these
trade-offs explicit and test whether degraded resource modes preserve calibrated
abstention rather than confident error.

\section*{Supplementary Note 13: Minimum cross-layer reporting standard}
A complete report should identify: the model and prompt; context and retrieval
policy; memory write/read/forget rules; tool schemas and error semantics;
protocol and authorization assumptions; harness orchestration and recovery;
situation schema, update rules and action gate; evaluation data, judges and
confidence intervals; observability and provenance records; and serving budgets,
latency, cost and failure modes. Reporting only the model name and prompt is
insufficient for agentic systems because performance is jointly determined by
the model, environment and scaffolding.

\end{document}
